ComplianceAI is built as a three-layer system: a document processing layer
that builds and maintains the regulatory knowledge base, an inference backend
that handles retrieval and LLM-based reasoning, and a web interface that lets
users submit documents and read compliance reports.
Figure~\ref{fig:architecture} provides an overview of the full pipeline.

\subsection{Knowledge Base Layer}

The knowledge base consists of nine official YÖK regulatory documents,
covering topics from graduate admission requirements to financial aid rules
and international student procedures.
These documents were collected as PDFs from the YÖK website
(\url{https://www.yok.gov.tr}) and from Istanbul Bilgi University's
official resources.

\paragraph{Current approach --- manual PDF management.}
At this stage, the regulatory documents are maintained manually: when YÖK
publishes an update, the corresponding PDF is replaced in the local
document store and the indexes are rebuilt.
This works well for a prototype but is not ideal for production.

\paragraph{Future direction --- live API integration.}
YÖK publishes regulations through its official web portal, and future
versions of ComplianceAI could automate document updates by periodically
fetching new or revised regulations directly from the source.
This would eliminate the manual maintenance step and ensure the system
always reflects the current regulatory landscape --- particularly important
as YÖK regulations are updated several times per year.

Each document is split into overlapping text chunks using a sliding window:
350 words per chunk with a 59-word overlap.
The resulting corpus contains 248 chunks across nine source documents.
Two index structures are built from these chunks.
For sparse retrieval, we build a BM25 index using \texttt{rank\_bm25}.
For dense retrieval, we encode each chunk with OpenAI's
\texttt{text-embedding-3-small} model (projected to 384 dimensions)
and store the vectors in a FAISS flat inner-product index.
Both indexes are persisted to disk so they do not need to be rebuilt
on each request.

\subsection{Inference Backend}

The inference backend is implemented as a FastAPI application.
When a user submits a document and a compliance question, the backend
routes the request to one of five retrieval pipelines:

\begin{itemize}
  \item \textbf{No-RAG}: The document and question go directly to the LLM
    without any retrieved context. This is the baseline.
  \item \textbf{BM25 (Sparse)}: The query is scored against the BM25 index
    and the top-$K$ chunks are retrieved.
  \item \textbf{Dense}: The query is embedded and compared against the FAISS
    index. The top-$K$ most similar chunks are retrieved.
  \item \textbf{Hybrid}: BM25 and Dense are run independently.
    Results are merged using reciprocal rank fusion (equal 0.5/0.5 weighting)
    and the top-$K$ are kept.
  \item \textbf{Multi-Query}: The original query is expanded into three
    related sub-queries via a lightweight LLM call.
    Each sub-query is run through BM25 and results are merged.
\end{itemize}

The retrieved chunks are formatted into a structured prompt and sent to
GPT-4o or GPT-4o-mini via the OpenAI API.
The system prompt instructs the model to respond only with a JSON object
containing four fields: \texttt{compliance\_score} (0--100),
\texttt{label} (compliant / partial / non-compliant),
\texttt{explanation} (a plain-language Turkish summary), and
\texttt{articles} (a list of specific YÖK articles referenced).
Label thresholds: compliant $\geq 70$, partial $30$--$69$,
non-compliant $< 30$.
The full system prompt is listed in Appendix~\ref{app:prompt}.

\subsection{Web Interface}

\textit{[This section will be updated with final implementation details
once the frontend and backend deployment are complete.
The current prototype includes a React-based dashboard and a FastAPI
REST API. Screenshots and endpoint documentation will be added here.]}

The interface is designed to be usable by non-technical administrative staff.
Users can upload a policy document, enter a compliance question, select the
retrieval pipeline, and receive a structured compliance report in seconds.
Results are displayed with a color-coded compliance score gauge, a
plain-language explanation, and the specific YÖK articles that were
retrieved and used in the assessment.

\subsection{Generalizability}

Although ComplianceAI was built and evaluated using Istanbul Bilgi University
documents and YÖK regulations, the system architecture is not institution-specific.
Any Turkish university could use the same system by loading their own
institutional documents and the same YÖK regulatory knowledge base.
Beyond Turkey, the same RAG-based approach could be adapted for compliance
checking in other higher education regulatory systems --- for example,
ECTS regulations in Europe or ABET accreditation requirements in the US ---
by simply replacing the regulatory document collection.
The benchmark construction methodology we describe in Section~\ref{sec:methodology}
could also be reused to evaluate compliance systems in other domains.
